\chapter{Introduction}
\label{ch:intro}

Software maintenance is a continuous and expensive part of the software life cycle. After initial development, software systems must be corrected, adapted, improved, and extended to keep them operational, reliable, and aligned with changing requirements, platforms, dependencies, and user expectations~\cite{kafura_use_1987, borstler_role_2016, sokolowski2022change, zhang2024developers}. This maintenance effort applies not only to production code but also to test code. Since test code is used to check the correctness of production code, its own correctness and maintainability are equally important. If tests are incorrect, outdated, or difficult to maintain, they may fail to detect regressions, report false failures, or provide developers with misleading confidence. Therefore, test code should not be treated as a secondary artifact; it deserves research and maintenance attention comparable to production code.

To support software maintenance, researchers have developed indicators and tools that help identify artifacts requiring developer attention, such as source-code metrics, code smells, change histories, and self-admitted technical debt~\cite{romano_using_2011, khomh_exploratory_2009, pascarella_performance_2020, chowdhury_method-level_2024, chowdhury_good_2025, potdar_exploratory_2014, maldonado_detecting_2015}. In this thesis, we examine two such indicators in the context of test code: self-admitted technical debt, which captures maintenance concerns explicitly documented by developers~\cite{potdar_exploratory_2014, maldonado_detecting_2015}, and method-level change history, which shows how individual methods evolve over time~\cite{higo_tracking_2020, grund_codeshovel_2021, jodavi_accurate_2022}. However, both lines of research have focused mainly on production code.

This production-code focus leaves important questions about test code unanswered. Although prior studies have shown that test code has its own quality and maintenance concerns, including test smells and cloning~\cite{bavota_are_2015, tufano_empirical_2016, counsell_comparing_2016, van_bladel_comparative_2023, spadini_relation_2018}, we still know less about how self-admitted technical debt appears in test code and how test methods evolve over time.

This thesis addresses this gap by studying test-code maintenance from two complementary perspectives. First, it explores self-admitted technical debt (SATD) in test code, where developers explicitly record maintenance concerns, shortcuts, workarounds, or future work in comments. Second, it investigates method-history generation and method-level test-code evolution, focusing on how individual methods can be tracked accurately across revisions and how test methods evolve relative to the production methods they exercise. Together, these two perspectives provide a more detailed understanding of test-code maintenance and help build stronger empirical foundations for future research on test-code quality and evolution.

\section{Background and Motivation}

\subsection{Technical Debt}

Technical debt refers to suboptimal implementation decisions that may provide short-term benefits but increase future maintenance costs~\cite{seaman_technical_2015}. A particular form of technical debt is self-admitted technical debt (SATD), where developers explicitly document debt, shortcuts, workarounds, or postponed work in code comments~\cite{potdar_exploratory_2014}. Potdar and Shihab~\cite{potdar_exploratory_2014} were among the first to systematically study SATD in open-source projects, and Maldonado and Shihab~\cite{maldonado_detecting_2015} later classified SATD comments into categories such as design, requirement, defect, test, and documentation debt. Subsequent studies showed that SATD can persist over time~\cite{bavota_large-scale_2016}, negatively affect software evolvability~\cite{wehaibi_examining_2016}, and be associated with larger, more complex, less readable, more change-prone, and more bug-prone methods~\cite{chowdhury_evidence_2025}.

Because SATD captures maintenance concerns explicitly written by developers, it has received considerable research attention. Prior studies have proposed techniques for detecting SATD~\cite{maldonado_detecting_2015, liu_satd_2018, huang_identifying_2018, li_automatic_2023, sheikhaei_empirical_2024, maldonado_using_2017, gama_towards_2024, rantala_prevalence_2020}, classifying SATD~\cite{zhu_detecting_2021, maldonado_detecting_2015, obrien_artifact_2022, bhatia_empirical_2025}, and mitigating or tracking SATD over time~\cite{maldonado_empirical_2017, zampetti_was_2018, zampetti_automatically_2020, iammarino_empirical_2021, liu_exploratory_2021, mastropaolo_towards_2023}. Researchers have also extended SATD studies to specialized domains, including machine learning software~\cite{obrien_artifact_2022, bhatia_empirical_2025}, Dockerfiles~\cite{azuma_empirical_2022}, blockchain projects~\cite{pinna_investigation_2023}, and Android applications~\cite{wilder_exploratory_2023}. However, the subject systems and datasets used in most SATD studies have been either exclusively production code or overwhelmingly dominated by production code. Although Counsell and Swift~\cite{counsell_timing_2022} examined SATD in both test and production classes, their focus was the timing of SATD introduction rather than the characterization or detection of SATD in test code.

Relying on findings from production-code SATD to understand test-code SATD may lead to incomplete or inaccurate conclusions. Test code differs from production code in several important ways. Prior studies have shown that test code can contain test-specific code smells~\cite{bavota_are_2015, tufano_empirical_2016}, exhibit a higher prevalence of cloning~\cite{van_bladel_comparative_2023}, and contain comments less frequently than production code~\cite{counsell_comparing_2016}. Moreover, developers may devote less attention to the quality of test code than to production code, which can result in obsolete, fragile, difficult-to-maintain, or low-quality tests~\cite{counsell_comparing_2016, beller_developer_2019, beller_when_2015, zaidman_studying_2011, athanasiou_test_2014, spadini_relation_2018}. These differences point to the possibility that SATD in test code may have different forms, causes, and detection challenges than SATD in production code.

Therefore, SATD in test code remains insufficiently understood. It is still unclear whether SATD in test code follows the same patterns as SATD in production code, whether production-oriented SATD taxonomies are sufficient for test code, and whether existing SATD detectors or large language models can reliably identify SATD in test comments. This motivates Chapter~\ref{ch:td}, which characterizes SATD in test code and evaluates the effectiveness of existing tools and large language models for detecting it.

\subsection{Method Evolution}
Software evolution histories provide valuable evidence for understanding how software systems change over time. Practitioners use code histories for provenance analysis, traceability, code comprehension, onboarding, code review, and expert identification~\cite{grund_codeshovel_2021}. Researchers use them to build maintenance models for bug prediction, change prediction, clone management, technical debt analysis, authorship attribution, and maintenance effort estimation~\cite{duala-ekoko_tracking_2007, kashiwabara_method_2015, meng_mining_2013, rahman_ownership_2011, pascarella_performance_2020, chowdhury_method-level_2024, chowdhury_good_2025, chowdhury_evidence_2025}. However, version control systems such as \texttt{Git} primarily record line-level changes, whereas many maintenance questions require finer granularity. Prior work has proposed tools for fine-grained (method-level) history reconstruction. \textit{Historage}~\cite{hata_historage_2011} and \textit{FinerGit}~\cite{higo_tracking_2020} provide method-level analysis by restructuring repositories, but they require substantial preprocessing. \textit{CodeShovel} supports on-demand method-history extraction using lightweight similarity-based heuristics~\cite{grund_codeshovel_2021}. \textit{CodeTracker} later improved method-history reconstruction by correcting part of the \textit{CodeShovel} oracle and incorporating refactoring aware tool~\cite{jodavi_accurate_2022}. Despite these advances, method-level history reconstruction remains difficult because methods can be renamed, moved, refactored, or substantially modified across commits~\cite{higo_tracking_2020, grund_codeshovel_2021, jodavi_accurate_2022}. Tool evaluations also depend on reliable oracles, but manually reconstructing complete method histories is error-prone because annotators must inspect long sequences of commits. Inaccurate oracles can affect conclusions about tool accuracy and downstream studies that rely on method histories. This motivates Chapter~\ref{ch:hf}, which constructs systematically validated method-history oracles and introduces \textit{HistoryFinder}, a method-level history generation tool designed to improve accuracy while maintaining competitive runtime.

Reliable method histories are also necessary for studying test-method evolution, where prior work has examined production-test co-evolution, test-suite changes, test smells, and obsolete or deleted tests~\cite{zaidman2011studying, wang2021understanding, sun_revisiting_2023, pinto2012understanding, Spadini:2018, hao2013bug, bhatta2025understanding, yaraghi2025automated}. However, many studies are constrained by the lack of validated test method histories and rely on custom heuristics or extensive manual inspection~\cite{pinto2012understanding, bhatta2025understanding, Spadini:2018}. Moreover, recent method-history tools were mainly developed and evaluated using production-code oracles, so their effectiveness on test methods remains unclear. However, tracking test-method histories alone is not enough to explain why test methods evolve, since changes may be driven by production-code changes or by test-specific issues such as obsolete assertions, test smells. Distinguishing between these possibilities requires method-level links between test methods and the production methods they exercise. Existing production-to-test mapping techniques employ both static and dynamic approaches~\cite{van_rompaey_establishing_2009, white_establishing_2020, white_tctracer_2022, sun_method-level_2024}, but they have often been evaluated using limited ground truth from a limited number of projects and can involve substantial precision--recall trade-offs. These limitations motivate carefully assessing mapping approaches before using them in downstream studies of test-method evolution.




Taken together, these limitations reveal three main research gaps. First, existing SATD research provides limited evidence about how developers admit technical debt in test code. Second, method-level history generation still suffers from oracle and tooling limitations. Third, test-code evolution has not been sufficiently studied using validated test method histories and method-level production-to-test links. Without these foundations, it is difficult to determine whether test methods mainly evolve in response to production code or due to their own maintenance concerns. Therefore, the objective of this thesis is to advance understanding of test code by investigating its technical debt and method-level evolution. Specifically, this thesis studies SATD in test code, develops and evaluates accurate method-history generation support, and uses validated histories and production-to-test links to analyze how test methods evolve relative to the production methods they exercise. The thesis focuses on open-source Java systems and uses method-level analysis as the primary granularity because methods are small enough to capture fine-grained maintenance behavior but large enough to represent meaningful units of developer work.

\section{Thesis Contributions}

This thesis makes three main contributions, presented in Chapters~\ref{ch:td}, \ref{ch:hf}, and~\ref{ch:t2p}. The first contribution corresponds to Chapter~\ref{ch:td}, which examines self-admitted technical debt in test code. The second and third contributions correspond to Chapters~\ref{ch:hf} and~\ref{ch:t2p}, which first establish accurate method-history generation support and then use that support to examine method-level test-code evolution.

\subsection{Self-Admitted Technical Debt in Test Code}

The first contribution is a large-scale empirical study of SATD in test code. Chapter~\ref{ch:td} manually analyzes 50,000 test-code comments sampled from 1.6 million comments across 1,000 open-source Java projects. SATD comments in test code are identified, classified into a taxonomy that includes test-specific categories, and used to evaluate existing SATD detection tools as well as open-source and proprietary large language models.

The results show that test-code SATD has characteristics that are not fully captured by production-code SATD taxonomies or detectors. Existing SATD tools can provide useful precision but limited recall, while the evaluated large language models struggle to provide reliable precision-recall trade-offs for this task. This contribution provides a manually curated dataset, a test-code SATD taxonomy, and empirical evidence that can support future research on test-aware SATD detection and classification.

\subsection{Accurate Method History Generation}

The second contribution is a methodological foundation for accurate method-level evolution research. Chapter~\ref{ch:hf} constructs systematically developed method-history oracles using a semi-automated process that combines multiple state-of-the-art tools with expert manual validation. It also introduces \textit{HistoryFinder}, a method-level history generation tool designed to improve accuracy while maintaining competitive runtime performance.

The evaluation shows that \textit{HistoryFinder} provides a strong overall balance of precision, recall, F1 score, and runtime among the evaluated research tools. This contribution improves both the benchmark data and the tool support needed for method-level maintenance research on change-proneness, bug-proneness, technical debt, and test-code evolution.

\subsection{Test Method Evolution}

The third contribution is a method-level empirical study of test-code evolution. Chapter~\ref{ch:t2p} first evaluates whether state-of-the-art method-history tools can accurately track test method histories using a manually constructed oracle. It then evaluates method-level production-to-test mapping approaches utilising existing benchmarks and a newly constructed oracle. Using the most suitable history-tracking and mapping approaches, the revision frequencies of test methods and their corresponding production methods are compared, and the association between test smells and highly revised test methods is examined.

The results challenge the common assumption that test code primarily changes in response to production code. A substantial portion of test methods undergo as many or more revisions than their corresponding production methods, and highly revised small test methods show stronger associations with test smells. This contribution provides a more precise empirical basis for future research on test maintenance, regression testing, obsolete test detection, test repair, and test smell management.

Overall, this thesis makes the following contributions:

\begin{itemize}
\item It provides a manually curated dataset and taxonomy of SATD in test code and evaluates the ability of existing SATD tools and large language models to detect test-code SATD.
\item It constructs reliable method-level history oracles and introduces \textit{HistoryFinder}, a method-history generation tool with strong accuracy and competitive runtime.
\item It evaluates method-history tools on test methods and evaluates method-level production-to-test mapping techniques using manually constructed benchmarks.
\item It presents a method-level empirical study showing how test methods evolve relative to the production methods they exercise and how test smells are associated with highly revised test methods.
\end{itemize}

\section{Thesis Organization}

The remainder of this thesis is organized into two parts after this introductory chapter.

Part~\ref{part:technical-debt}, Self-Admitted Technical Debt in Test Code, contains Chapter~\ref{ch:td}, which examines how developers admit technical debt in test code and how effectively such debt can be detected. The chapter describes the collection and manual analysis of test-code comments, develops a taxonomy of test-code SATD, and evaluates existing SATD detectors and large language models for identifying SATD in test comments.

Part~\ref{part:test-method-evolution}, Method-History Generation and Test Method Evolution, contains Chapters~\ref{ch:hf} and~\ref{ch:t2p}. This part develops the methodological support needed for method-level evolution research and then applies that support to study test methods. Chapter~\ref{ch:hf} discusses limitations in existing method-history oracles, describes the construction of new oracles, introduces \textit{HistoryFinder}, and evaluates method-history tools. Chapter~\ref{ch:t2p} then presents the study of method-level test-code evolution by evaluating method-history tools on test methods, evaluating production-to-test method linking approaches, comparing the revision frequencies of test and production methods, and analyzing the association between test smells and highly revised test methods.

Chapter~\ref{ch:conclusion} concludes the thesis by summarizing the findings from both parts, discussing their implications for test-code maintenance research, and outlining directions for future work.
